\documentclass[a4paper,11pt]{article}
\pdfoutput=1

\usepackage{jinstpub}
\usepackage{booktabs}
\usepackage{lineno}
\usepackage{microtype}
\renewcommand{\ttdefault}{cmtt}

\graphicspath{{figures/}}
\linenumbers

\title{\boldmath From single-fiber occupancy to full-detector response:
digital SiPM requirements for HG-DREAM and FLASH}

% TODO before circulation: confirm the complete author list, ordering,
% affiliations, corresponding author, and email address.
\author[a]{C. Madrid}
\affiliation[a]{Department of Physics and Astronomy, Texas Tech University,\\
Lubbock, Texas 79409, U.S.A.}

\abstract{FLASH is a digital silicon-photomultiplier concept that replaces
integrated waveforms with sparse photon counts and timestamps for
time-resolved optical calorimetry. We report the next stage of a program that
began with a one-fiber HG-DREAM occupancy study and now supports a funded SLAC
LDRD prototype effort. An exact reproduction of the earlier analysis revealed
that its 3 mm segmentation had been applied to a 1 mm fiber-local display
coordinate, so the labeled pitch acted as a three-times finer physical pitch.
We retain that result as software closure and repeat the calculation on a
centered 3$\times$3 mm$^2$ sensor in physical millimeters. The corrected study
uses 5,000 simulated 100 GeV positron events and 5,000 simulated 1 TeV
charged-pion events to scan 5--100 $\mu$m pitch, 10--100\% constant
photon-detection efficiency, and 1, 4.3, and 40 ns recovery. At 20\% PDE and
4.3 ns recovery, 10, 20, and 25 $\mu$m cells retain 97.93\%, 92.73\%, and
89.50\% of detected photons in the pion stress sample. The corresponding
10th-percentile event values are 96.7\%, 90.3\%, and 86.2\%. A simple
longitudinal timing proxy remains limited by shower and optical-transport
tails rather than sub-100 ps detector smearing. We then propagate the sensor
response through available full-detector Cherenkov samples. After electron
calibration, a representative conventional C-Series SiPM proxy produces a
+9.47\% pion Cherenkov-response shift, compared with +1.67\% for a 20
$\mu$m, 4.3 ns FLASH point. The comparison establishes a reproducible chain
from Geant4 photons to detector-level response and converts the prototype
program into measurable requirements for PDE, recovery, timing response, and
output policy. Matched-energy samples with explicit Cherenkov and scintillation
photons remain necessary for a final calorimeter-resolution claim.}

\keywords{Calorimeters; photon detectors; silicon photomultipliers; timing detectors; detector simulation}

\begin{document}
\maketitle
\flushbottom

\section{Introduction}
\label{sec:introduction}

Digital silicon photomultipliers (dSiPMs) combine Geiger-mode avalanche
photodiodes with local quenching, counting, and timing logic
\cite{frach2009dsipm}. In a calorimeter, such a sensor could represent the
optical signal as sparse hit records rather than a continuously sampled analog
waveform. The FLASH concept targets direct single-photon timestamps at the
order of 10--50 ps and seeks to preserve spatial and temporal information from
the optical shower.

The prototype program uses working engineering targets of no more than about
20 $\mu$m pitch, at least 70\% active-area fill factor, and approximately
50 ps timing. These are development targets rather than performance values
established by the present calorimeter study. In particular, physical pitch,
photosensitive junction size, fill factor, photon-detection probability, and
the size of the shared timing channel are distinct quantities.

HG-DREAM provides a useful benchmark for this development. It combines dense
Cherenkov and scintillating fiber readout with fine transverse segmentation
and a 2 m optical propagation distance \cite{akchurin2025hgdream,feng2026hgdream}.
The sensor must tolerate bright, localized photon densities, while the fiber
length provides a timing lever arm for longitudinal shower information. The
same system also has conventional SiPM energy and waveform readout, providing
a path to simulation-to-data closure.

The CPAD 2025 study introduced this approach through a one-fiber HG-DREAM
simulation, reporting occupancy trends versus SPAD size and dead time together
with an initial timing-based reconstruction study~\cite{wu2025cpad}. The
present work is the next stage of that program and provides the simulation
foundation for a funded SLAC Laboratory Directed Research and Development
effort toward a FLASH prototype.

We first reproduce the earlier occupancy study exactly. The reproduction
closes to the saved ROOT histograms but exposes a coordinate-convention error:
the old labeled pitch was not the physical cell width in the DREAMSim
coordinate system. This distinction materially changes the design conclusion.
We therefore separate four products throughout the paper:
\begin{enumerate}
  \item historical closure, which establishes software provenance;
  \item a corrected physical sensor scan, which constrains pitch, PDE, and
  recovery jointly; and
  \item a first timing-information test, which identifies the next simulation
  and prototype measurements; and
  \item a full-detector Cherenkov response comparison between a conventional
  SiPM proxy and FLASH.
\end{enumerate}

The two available large samples have different beam energies. We use them as
occupancy stress cases and do not interpret their difference as a controlled
electron--pion comparison.

\section{Simulation samples and input audit}
\label{sec:samples}

The optical photons were produced with an HG-DREAM-like Geant4 simulation
\cite{geant4}. The analysis selects Cherenkov photons in the fiber core,
produced in fiber index zero, exiting the downstream end, and arriving in the
open interval $0<t<40$ ns. The historical fiber-center translation is applied,
but the resulting detector-plane coordinates remain in physical millimeters.
The corrected sensor covers $-1.5\leq x,y<1.5$ mm.

Table~\ref{tab:cutflow} reports the sequential cutflow. In these files every
stored photon already satisfies the valid-fiber, Cherenkov-core, and
downstream-exit selections. The time gate removes 4,321 photons from the
positron file and 13,543 photons from the pion file. The centered 3 mm sensor
contains every selected fiber-zero photon. The selected coordinates span
approximately $-0.39$ to $+0.40$ mm in each direction, consistent with the
face of one simulated 0.8 mm fiber.

\begin{table}[htbp]
\centering
\caption{\label{tab:cutflow}Sequential input audit for the two occupancy
stress samples.}
\smallskip
\begin{tabular}{lrr}
\toprule
Selection & 100 GeV $e^+$ & 1 TeV $\pi^+$ \\
\midrule
Events & 5,000 & 5,000 \\
Raw optical photons & 2,629,165 & 7,722,345 \\
Inside 0--40 ns & 2,624,844 & 7,708,802 \\
Production fiber zero & 649,310 & 1,899,950 \\
Inside centered 3 mm sensor & 649,310 & 1,899,950 \\
Median selected photons/event & 127 & 289 \\
90th percentile photons/event & 173 & 920 \\
99th percentile photons/event & 222 & 1,484 \\
\bottomrule
\end{tabular}
\end{table}

Figure~\ref{fig:xy} shows the two detector-plane photon distributions on a
common scale. The pion file has a wider occupancy distribution and reaches a
maximum of 2,108 selected photons in one fiber. Its higher energy and larger
event tails make it the more restrictive of the two available stress samples.

\begin{figure}[htbp]
\centering
\includegraphics[width=0.98\textwidth]{photon_xy_comparison.pdf}
\caption{\label{fig:xy}Selected photon density at the simulated fiber face for
the 100 GeV positron sample (left) and 1 TeV pion sample (right). Both panels
average 5,000 events and use the same color scale. The files have different
beam energies and are not a particle-species comparison.}
\end{figure}

\subsection{Coordinate-mapping audit}
\label{sec:mapping}

The historical study set the number of bins on one axis to
\begin{equation}
  N_{\mathrm{bin}}=\frac{3000~\mu\mathrm{m}}{p_{\mathrm{label}}},
\end{equation}
where $p_{\mathrm{label}}$ was the pitch printed on the plot. Those bins were
then applied across a fiber-local coordinate range of only 1 mm,
$[-0.5,+0.5)$ mm. The implemented cell width was therefore
\begin{equation}
  p_{\mathrm{implemented}}=
  \frac{1~\mathrm{mm}}{N_{\mathrm{bin}}}
  =\frac{p_{\mathrm{label}}}{3}.
  \label{eq:scale}
\end{equation}
A labelled 25 $\mu$m point behaved like an approximately 8.3 $\mu$m cell in
physical DREAMSim coordinates.

The revised calculation leaves photon coordinates in millimeters and assigns
cells directly over the centered 3$\times$3 mm$^2$ sensor. Figure
\ref{fig:mapping} quantifies the effect. At unity PDE, 25 $\mu$m pitch, and
4.3 ns recovery, the historical pion result retains 93.24\% of detected
photons, while the physical mapping retains 65.18\%. The corresponding
positron values are 98.73\% and 90.51\%. Exact closure of the blue curves is
still useful, but those curves cannot define a physical FLASH pitch.

\begin{figure}[htbp]
\centering
\includegraphics[width=0.98\textwidth]{v2_mapping_correction.pdf}
\caption{\label{fig:mapping}Recovery survival versus labelled pitch with the
historical display-scaled mapping and with physical millimeter coordinates.
Both calculations use unity PDE, 4.3 ns recovery, and one Cherenkov fiber.}
\end{figure}

\section{Sensor-response model}
\label{sec:model}

For each event, selected photons are sorted by true arrival time. A constant
PDE value $\epsilon$ is applied with a deterministic per-photon random number,
so all scan points use common random numbers. Detected photons are assigned to
square cells of physical pitch $p$. Photons outside the sensor are rejected
rather than clipped to edge cells.

The recovery model is non-paralyzable. A detected photon $i$ in cell $c$ is
accepted when
\begin{equation}
  t_{i,c}-t^{\mathrm{last}}_c\geq \tau_{\mathrm{rec}},
  \label{eq:recovery}
\end{equation}
where $t^{\mathrm{last}}_c$ is the time of the most recent accepted avalanche.
A rejected photon does not extend the holdoff. The scan uses pitches of 5, 10,
15, 20, 25, 30, 40, 50, 75, and 100 $\mu$m; PDE values of 10\%, 20\%, 30\%,
50\%, and 100\%; and recovery times of 1, 4.3, and 40 ns. The 4.3 ns point
matches a public Litavis holdoff specification \cite{litavis2026}. The 40 ns
point approximates one accepted hit per cell in the event gate.

We report two ratios:
\begin{align}
  U &= \frac{N_{\mathrm{accepted}}}{N_{\mathrm{incident}}},
  &
  S &= \frac{N_{\mathrm{accepted}}}{N_{\mathrm{detected}}}.
  \label{eq:metrics}
\end{align}
The usable fraction $U$ includes PDE loss. Recovery survival $S$ isolates
repeated-cell loss after detection. Reporting only $S$ can make a low-PDE
sensor appear artificially robust because photon rejection relieves cell
collisions. Aggregate confidence intervals use 500 bootstrap resamples of
complete Geant4 events.

The recovery and cell-assignment code is independent of ROOT input/output and
has unit tests for out-of-sensor rejection, finite recovery, event-long
holdoff, common random numbers, timestamp ordering after jitter, grouping, and
finite-cell saturation. The saved historical electron histograms close
exactly in binning and integral, and an independent 40 ns implementation
matches the historical one-hit calculation at every shared scan point.

\section{Physical dynamic-range results}
\label{sec:results}

Figure~\ref{fig:requirements} shows the physical requirement surface for the
1 TeV pion stress sample. At fixed 20\% PDE, survival degrades rapidly as
pitch and recovery time increase. At fixed 4.3 ns recovery, increasing PDE
also increases collision loss. A pitch requirement therefore cannot be quoted
without an operating PDE and recovery definition.

\begin{figure}[htbp]
\centering
\includegraphics[width=0.98\textwidth]{v2_physical_requirements.pdf}
\caption{\label{fig:requirements}Recovery survival in the 1 TeV pion stress
sample. Left: physical pitch for three recovery times at 20\% PDE. Right: PDE
for three physical pitches at 4.3 ns recovery. Dashed horizontal lines mark
90\% and 95\% survival.}
\end{figure}

Table~\ref{tab:working} gives values in the region most relevant to current
fine-pitch architectures. At 20\% PDE and 4.3 ns recovery, 10 $\mu$m pitch
loses only 2.07\% of detected photons to recovery. The loss grows to 7.27\%
at 20 $\mu$m and 10.50\% at 25 $\mu$m. The usable fraction changes by less
than two percentage points across these rows because the assumed 20\% PDE is
the leading absolute-yield term.

\begin{table}[htbp]
\centering
\caption{\label{tab:working}Physical-pitch results for the 1 TeV pion stress
sample at 20\% constant PDE and 4.3 ns recovery.}
\smallskip
\begin{tabular}{rccc}
\toprule
Pitch & Usable fraction $U$ & Recovery survival $S$ & 10th-percentile event $S$ \\
\midrule
10 $\mu$m & 19.64\% & 97.93\% & 96.7\% \\
15 $\mu$m & 19.16\% & 95.57\% & --- \\
20 $\mu$m & 18.59\% & 92.73\% & 90.3\% \\
25 $\mu$m & 17.95\% & 89.50\% & 86.2\% \\
30 $\mu$m & 17.23\% & 85.92\% & --- \\
\bottomrule
\end{tabular}
\end{table}

The aggregate ratios are photon weighted and therefore emphasize bright
events. They still do not expose how performance varies between events.
Figure~\ref{fig:tails} shows the cumulative event-level survival at 20\% PDE
and 4.3 ns recovery. At 20 $\mu$m, aggregate survival is 92.73\%, but the
10th-percentile event survives at only 90.3\%. At 25 $\mu$m the corresponding
values are 89.50\% and 86.2\%. A device requirement should specify both an
aggregate metric and an allowed bright-event tail.

\begin{figure}[htbp]
\centering
\includegraphics[width=0.82\textwidth]{v2_event_tail.pdf}
\caption{\label{fig:tails}Cumulative event-level recovery survival in the
1 TeV pion stress sample at 20\% PDE and 4.3 ns recovery. Events with no
PDE-selected photons are excluded from the ratio.}
\end{figure}

Repeated-cell loss is also time dependent. Figure~\ref{fig:pulse} begins with
the same PDE-selected photon stream and compares a conventional analog SiPM
partial-recharge model with several dSiPM pitch and holdoff scenarios.  The
dSiPM hit streams are convolved with one common display kernel so that they can
be read using the familiar visual language of an oscilloscope.  These curves
are waveform-equivalent renderings, not predicted dSiPM output voltages.  The
lower panel retains the absolute cumulative normalization and therefore shows
both pulse-shape distortion and total loss.

\begin{figure}[htbp]
\centering
\includegraphics[width=0.96\textwidth]{pulse_architecture_comparison.pdf}
\caption{\label{fig:pulse}Mean one-fiber pulse response for the 1 TeV pion
stress sample using a common 20\% PDE selection.  The conventional-SiPM curve
uses exponential partial recharge for a representative 3 mm Onsemi sensor.
The digital curves use hard holdoff at the stated SPAD pitches and are rendered
with the same 0.6 ns rise, 1.5 ns FWHM display kernel.  Public-device PDE,
timing jitter, grouping, and exported timestamp policy are excluded here to
isolate pitch and recovery.  The 200 $\mu$m curve is a simulation stress point,
not a fabricated-device claim; the 40 ns holdoff paired with the Philips pitch
is likewise a comparison proxy, not a DPC3200 specification.}
\end{figure}

Figure~\ref{fig:pulse-scan} extends the scan to 200 $\mu$m and separates
integrated response loss from the shift of the mean photon time.  The 78
$\mu$m Sherbrooke PDC is the largest pitch of a fabricated dSiPM-class array
identified in the reviewed primary material; 100--200 $\mu$m points are kept
as explicit design stress tests.  The comparison also demonstrates why SPAD
pitch and recovery must remain independent requirements.

\begin{figure}[htbp]
\centering
\includegraphics[width=0.96\textwidth]{pulse_pitch_recovery_scan.pdf}
\caption{\label{fig:pulse-scan}Integrated output and mean-time bias versus
SPAD pitch for hard holdoffs from 1 to 40 ns.  The purple star is the
conventional-SiPM partial-recharge proxy at its effective square-grid pitch;
it is not mathematically equivalent to a hard dead time.}
\end{figure}

\section{Factorized active area, spectral response, and channel sharing}
\label{sec:factorized}

The constant-PDE scan isolates occupancy and recovery, but it cannot test the
simultaneous pitch and fill-factor target. We therefore add a factorized
requirement model in which a photon must first land in the centered active
region of its pitch cell, then satisfy a wavelength-dependent junction
detection probability, survive physical-SPAD recovery, and finally obtain a
record from shared timing electronics. The sequential exported-timestamp
yield is
\begin{equation}
  \frac{N_{\mathrm{timestamp}}}{N_{\mathrm{incident}}}
  = \frac{N_{\mathrm{active}}}{N_{\mathrm{incident}}}
    \frac{N_{\mathrm{detected}}}{N_{\mathrm{active}}}
    \frac{N_{\mathrm{accepted}}}{N_{\mathrm{detected}}}
    \frac{N_{\mathrm{timestamp}}}{N_{\mathrm{accepted}}}.
  \label{eq:factorized_yield}
\end{equation}
This identity prevents an apparent improvement in recovery survival at low
fill factor from being mistaken for improved usable yield.

For the present requirement study, the active region is a centered square
with nominal area fraction 30--100\%. The wavelength shape is the normalized
polynomial used in the earlier capstone workflow, scaled to a 30\% peak
junction-detection probability. It is a traceable proxy rather than a FLASH
measurement. Figure~\ref{fig:spectral} shows the stored Cherenkov wavelength
distribution and each response stage for the 20 $\mu$m, 70\% fill-factor
point. The incident-spectrum-averaged junction probability is 26.85\%.
Geometric acceptance reduces the effective detected fraction to 18.74\%; SPAD
recovery retains 93.06\% of detected photons. An approximately 80 $\mu$m
shared timing channel retains 96.30\% of accepted avalanches, giving 16.79
exported timestamps per 100 incident photons.

\begin{figure}[htbp]
\centering
\includegraphics[width=0.98\textwidth]{spectral_response_proxy.pdf}
\caption{\label{fig:spectral}Left: wavelength distribution of photons in the
1 TeV pion one-fiber stress sample. Right: factorized response at 20 $\mu$m
pitch and 70\% nominal square active area. The spectral junction response has
a 30\% peak and is a study proxy, not a measured FLASH PDE curve.}
\end{figure}

Figure~\ref{fig:fillfactor} extends the calculation across pitch and fill
factor. Lower fill factor reduces collision probability and can therefore
increase recovery survival, while both effective detection and end-to-end
timestamp yield decrease. Pitch optimization must consequently be performed
at a stated active geometry and junction-response curve.

\begin{figure}[htbp]
\centering
\includegraphics[width=0.98\textwidth]{fill_factor_pitch_matrix.pdf}
\caption{\label{fig:fillfactor}Joint pitch and nominal square-fill-factor scan
for the 1 TeV pion stress sample. The four panels separate geometric active
area, effective detection, physical-SPAD recovery, and end-to-end timestamp
yield. The model uses the spectral proxy of figure~\ref{fig:spectral}, 4.3 ns
recovery, and an approximately 80 $\mu$m readout-channel side.}
\end{figure}

The final factor is controlled by timing-channel topology rather than SPAD
pitch alone. A second scan fixes PDE to 20\%, recovery to 4.3 ns, channel
resolving time to 50 ps, timestamp depth to eight, and counter depth to eight
bits. It varies physical pitch and square groups from $1\times1$ through
$64\times64$ SPADs, including a $6\times6$ topology. We select the
lowest-payload point that retains at least 95\% of timestamps, limits events
with any timing-buffer overflow to 5\%, and retains at least 99\% of counts.
Across 5--50 $\mu$m pitch, the selected nominal timing-channel side lies
between 50 and 100 $\mu$m. At 20 $\mu$m pitch, the $4\times4$ point retains
96.10\% of timestamps, has at least one overflow in 1.42\% of events, and
exports 1.617 kbit/event on average.

\begin{figure}[htbp]
\centering
\includegraphics[width=0.98\textwidth]{channel_architecture_matrix.pdf}
\caption{\label{fig:channels}Common-resource architecture stress test for the
1 TeV pion sample. Every point uses the same PDE, recovery, resolving time,
counter, and timestamp buffer. The result compares grouping under a common
output model; it is not a prediction of the native Litavis, DIGILOG, or
ASPIDES readout policies.}
\end{figure}

\section{First test of longitudinal timing information}
\label{sec:timing}

For a direct optical path, the propagation contribution is approximately
\begin{equation}
  t_{\mathrm{prop}}=\frac{nL}{c},
\end{equation}
where $L$ is the distance from photon production to the sensor and $n$ is the
fused-silica refractive index. Internal reflections and non-axial paths
produce a late tail. The time of an individual photon is therefore correlated
with production depth but does not determine it uniquely.

As a first information test, we use the 10th percentile of accepted photon
arrival times in each pion event. The front end is fixed to 20\% PDE, 20
$\mu$m pitch, and 4.3 ns recovery. Events require at least five accepted
photons and at least five selected photons with production depth between 0 and
2 m. A linear calibration predicts the mean selected-photon production depth.
Even-numbered events define the calibration and odd-numbered events form the
test set, giving 2,079 training and 2,033 test events.

The ideal-time test RMSE is 236.46 mm, but the median absolute error is only
59.24 mm and the 90th percentile absolute error is 246.77 mm. The large
difference indicates a small population of poorly modeled shower tails, so
this observable is a proof of information retention rather than a calorimeter
longitudinal-resolution measurement.

To isolate the incremental detector-timing contribution, each timing scenario
uses the same photons and Gaussian random variates. We refit the calibration
for each scenario and bootstrap the paired RMSE difference relative to ideal
timing. Figure~\ref{fig:timing} shows that 10--50 ps response/TDC terms have
negligible impact on this many-photon quantile. A 250 ps response plus 50 ps
TDC adds $2.24^{+1.11}_{-1.31}$ mm at the central 68\% paired interval. A
1.145 ns response plus 37 ps TDC adds 11.80 mm, with a broad interval from
1.19 to 30.15 mm. The result does not remove the motivation for 10--50 ps
single-photon timing, but it shows that this simple HG-DREAM observable does
not yet set that requirement.

\begin{figure}[htbp]
\centering
\includegraphics[width=0.84\textwidth]{v2_timing_depth.pdf}
\caption{\label{fig:timing}Added test RMSE of the pion depth proxy relative to
ideal photon times. Vertical bars give the central 68\% interval from a paired
event bootstrap. The ideal absolute RMSE is 236 mm and is dominated by the
proxy and shower tails rather than sub-100 ps detector smearing.}
\end{figure}

\section{First full-detector readout comparison}
\label{sec:full_detector}

We next apply the sensor response to every explicit Cherenkov photon in the
available full-detector files. This sample contains 1,000 100 GeV positron
events and 100 90 GeV charged-pion events. Each photon is assigned to its
physical rod, Cherenkov fiber, and local sensor cell before digitization. The
samples are not a matched-energy production and are used only to expose
particle-dependent response transfer.

Figure~\ref{fig:fulloccupancy} connects the one-fiber requirement to this
detector-level calculation. The brightest Cherenkov fiber in an electron
shower carries a larger local photon density than in the available pion
sample. The event-level retained fraction consequently depends on both the
sensor point and the spatial concentration of the shower.

\begin{figure}[htbp]
\centering
\includegraphics[width=0.98\textwidth]{full_detector_occupancy.pdf}
\caption{\label{fig:fulloccupancy}Local occupancy and retained event signal in
the exploratory full-detector samples. The left panels quantify the brightest
Cherenkov fiber in each event. The right panels compare the response after
recovery with the matched-PDE no-recovery baseline. Electron and pion samples
have different beam energies and statistics.}
\end{figure}

The conventional benchmark is a representative Onsemi MicroFC-30035 C-Series
proxy based on public data-sheet values: 31\% PDE at 420 nm and 2.5 V
overvoltage, 4,774 microcells, an 82 ns recharge constant, and a 1.5 ns
fast-output FWHM~\cite{onsemiMicroC}. A fully charged avalanche contributes
one photoelectron of charge. A repeated avalanche after $\Delta t$ contributes
$1-\exp(-\Delta t/82\,\mathrm{ns})$. The FLASH benchmark uses 20\% constant
PDE, 20 $\mu$m physical pitch, and a 4.3 ns non-paralyzable holdoff.

\begin{table}[htbp]
\centering
\caption{\label{tab:full_detector}Recovery response in the exploratory
full-detector samples. Each pion Cherenkov shift is measured after electron
calibration relative to a no-recovery baseline at the same PDE.}
\smallskip
\begin{tabular}{lccc}
\toprule
Readout point & Electron retained & Pion retained & Pion C shift \\
\midrule
Conventional C-Series proxy & 88.54\% & 96.93\% & +9.47\% \\
FLASH, 20 $\mu$m and 4.3 ns & 97.79\% & 99.43\% & +1.67\% \\
\bottomrule
\end{tabular}
\end{table}

The electron showers contain a much brighter hottest fiber than the pion
sample, so finite recovery compresses the electron calibration more strongly.
Applying that calibration to pions creates a positive response shift. The
smaller response width after either recovery model is therefore saturation
compression and must not be interpreted as improved calorimeter resolution.
The comparison isolates recovery by using a matched-PDE baseline for each
model; it does not claim equal photon yield for the 20\% and 31\% PDE points.

Figure~\ref{fig:fullscan} extends the two quoted benchmarks across the FLASH
pitch and recovery scan. The mean pion shift exposes response nonlinearity
after electron calibration. The width and compression panels diagnose how
saturation reshapes the distribution; a smaller width after saturation is not
by itself a resolution improvement.

\begin{figure}[htbp]
\centering
\includegraphics[width=0.98\textwidth]{full_detector_performance_scan.pdf}
\caption{\label{fig:fullscan}Full-detector response transfer for FLASH sensor
points and the representative conventional-SiPM proxy (purple star). Left:
pion Cherenkov-response shift after electron calibration, relative to a
matched-PDE baseline. Center: event RMS divided by the mean. Right: response
compression relative to that baseline.}
\end{figure}

This conventional-SiPM calculation is an optimistic sensor-only proxy, not a
model of the complete installed HG-DREAM chain. It omits wavelength-dependent
PDE, recharge-dependent avalanche probability, crosstalk, afterpulsing,
electronics noise, and waveform integration. Moreover, explicit optical
photons are available only for the Cherenkov channel. A final comparison must
digitize measured analog waveforms and FLASH sparse hits from a common,
matched-energy Cherenkov/scintillation photon production.

\section{Relation to current digital-SiPM architectures}
\label{sec:devices}

Table~\ref{tab:devices} places the corrected requirement beside selected
public architectures. The quantities are not interchangeable. Physical pitch
sets local occupancy, while timing-group size determines which photons share
electronics. A TDC least significant bit is not a complete single-photon
instrument response.

\begin{table}[htbp]
\centering
\caption{\label{tab:devices}Selected public benchmarks relevant to the FLASH
design space. Design targets and submitted architectures are distinguished
from measured full-system quantities in the cited sources.}
\smallskip
\begin{tabular}{p{2.1cm}p{2.6cm}p{2.7cm}p{5.0cm}}
\toprule
Architecture & Spatial scale & Timing quantity & Relevance \\
\midrule
FLASH & Architecture under study & 10--50 ps target & The present simulation supports a requirement surface, not a frozen physical pitch. \\
Litavis~\cite{litavis2026} & 10.17 $\mu$m SPAD; 40.68 $\mu$m macropixel & 37 ps TDC bin; 1,145 ps IRF at 450 nm & Demonstrates fine-pitch 3D integration and 4.3 ns holdoff; blue-light response is much wider than the TDC bin. \\
ASPIDES~\cite{aspides2025} & 30 $\mu$m SPAD & 100 ps TDC LSB; first/last timing & Calorimeter-oriented design with counting and compressed timing output. \\
DESY dSiPM~\cite{feindt2024desy} & about 70 $\mu$m pixel with four SPADs & $<100$ ps TDC; about 50 ps MIP result & Measured prototype; timing is shared within a larger pixel rather than exported per SPAD. \\
\bottomrule
\end{tabular}
\end{table}

The Litavis data sheet also reports a 47\% photon-detection efficiency at
785 nm, but it does not provide a numerical blue-light PDE suitable for the
present Cherenkov spectrum. We therefore do not insert that device into an
absolute HG-DREAM yield prediction. This is an example of why architecture,
wavelength response, and exported data semantics must be carried separately.

\section{Limitations and next studies}
\label{sec:limitations}

The corrected analysis begins with a local one-fiber study and then digitizes
every explicit Cherenkov photon in the available detector files. The stored
scintillation sum has no corresponding explicit photon population, so this is
not yet a complete dual-readout digitization. The calculation also omits
optical coupling, measured wavelength- and angle-dependent PDE, dark counts,
optical crosstalk, afterpulsing, hot-cell masking, time walk, non-Gaussian
timing tails, TDC nonlinearity, realistic arbitration, link bandwidth, and
power. The one-fiber architecture stress test includes idealized finite
counters and buffers, but these resources are not yet propagated through the
full-detector reconstruction. The constant-PDE, spectral-shape, active-area,
and partial-recharge models are requirement tools, not complete device
predictions.

Four developments now have higher value than additional scans of the current
files. First, matched electron, pion, and muon samples should be produced with
one DREAMSim configuration at several energies between approximately 10 and
300 GeV, retaining explicit Cherenkov and scintillation photons. This will
separate beam energy, particle type, shower density, and the two readout
channels.
Second, the sensor model should ingest measured FLASH and reference-device
PDE$(\lambda)$, instrument-response functions, and recovery distributions.
Third, the same physics endpoints should be reconstructed from all photon
timestamps, first/last timestamps, first hit per timing group, and short
histograms. The result should be reported versus bits per event and lost-hit
probability. Fourth, a full-tile simulation should be validated against
HG-DREAM beam-test waveforms before converting those photons into sparse hits.

Neural-network studies should follow, not precede, count-only and timing-only
baselines with frozen preprocessing, independent training and test samples,
and event-level uncertainty estimates. The first targets should be shower
depth, particle identification, and the improvement to a dual-readout energy
correction.

\section{Conclusions}
\label{sec:conclusions}

We reproduced an earlier HG-DREAM dSiPM occupancy analysis and used that
closure to find a consequential coordinate-convention error. A 3 mm sensor
segmentation had been applied over a 1 mm fiber-local display range, causing
the labelled pitch to act as a three-times finer physical pitch. Correcting
the mapping removes the earlier basis for a universal 25 $\mu$m FLASH working
point.

With physical millimeter coordinates, the 1 TeV pion stress sample at 20\%
PDE and 4.3 ns recovery retains 97.93\%, 92.73\%, and 89.50\% of detected
photons at 10, 20, and 25 $\mu$m pitch. Event-level tails and prompt-peak
distortion strengthen the constraint. The useful output of this phase is
therefore a PDE--pitch--recovery surface with an event-tail requirement.

A first photon-timing proxy confirms that the simulated arrival distribution
contains longitudinal information, but sub-100 ps response smearing does not
limit that simple many-photon observable. The next requirement must come from
matched-energy samples, measured optical response, and a comparison of
readout policies using calorimeter-level endpoints. HG-DREAM remains a strong
benchmark because it connects these sensor choices to dense optical showers,
current waveform readout, and beam data.

The first detector-level propagation already distinguishes the readout
concepts. Relative to matched-PDE no-recovery baselines, the representative
conventional-SiPM proxy retains 88.54\% of the electron charge and 96.93\% of
the pion charge, producing a +9.47\% calibrated pion Cherenkov shift. The 20
$\mu$m, 4.3 ns FLASH point retains 97.79\% and 99.43\%, respectively, and
reduces the shift to +1.67\%. The unmatched beam energies and absence of
explicit scintillation photons preclude a final calorimeter-resolution claim.

The factorized extension also shows why no single ``SPAD size'' specifies the
front end. At the illustrative 20 $\mu$m, 70\% fill-factor point, a 30\% peak
spectral junction proxy yields 18.74\% effective detection and 16.79\%
end-to-end timestamp yield. Under the separate common-resource architecture
scan, timing-channel sides of approximately 50--100 $\mu$m satisfy the stated
retention, overflow, and count criteria across the tested pitch range. These
values define testable design assumptions, not measured prototype
performance.

The funded SLAC LDRD now provides the experimental path to replace these
response proxies with measured FLASH inputs. Blue-light PDE, the
single-photon instrument response, recovery behavior, correlated noise, and
output policy can each be measured on the prototype and propagated to the
same one-fiber and detector-level observables. This closes the loop from
sensor development to calorimeter performance rather than treating the
prototype and simulation as separate efforts.

\section{Data and software availability}

The baseline repository is available at
\url{https://github.com/TTU-HEP/dSiPM}. The current draft includes a
deterministic response model, unit tests, input cutflows, scan tables, event
tables, timing-calibration results, and run manifests. A reviewed tagged
release with checksums and production metadata for the large derived Geant4
samples will be prepared before submission.

\acknowledgments

This work is supported in part by the SLAC Laboratory Directed Research and
Development program. The authors thank the HG-DREAM collaboration for the
detector simulation and continuing beam-test program. The complete author list
and funding acknowledgment will be confirmed before circulation.

Generative AI tools assisted code refactoring, validation visualization,
manuscript organization, and copy editing. The authors will verify all
analysis choices, numerical results, references, and scientific claims before
submission.

\bibliographystyle{JHEP}
\bibliography{references}

\end{document}
